% Fórmulas de um curso de Cálculo e Física. O Fullpad não compila documentos LaTeX:
% ele mostra, ao lado do código (Ctrl+E), cada equação solta do arquivo.
\documentclass[12pt]{article}
\usepackage{amsmath, amssymb}

\title{Notas de aula: limites, séries e mecânica}
\author{Turma 2026}

\begin{document}
\maketitle

\section{Limites e continuidade}

O limite fundamental é $\lim_{x \to 0} \frac{\sin x}{x} = 1$, e a definição
formal usa $\varepsilon$ e $\delta$:
\[
  \forall \varepsilon > 0,\ \exists \delta > 0 :\
  0 < |x - a| < \delta \implies |f(x) - L| < \varepsilon
\]

Uma função definida por partes:
\begin{equation}
  f(x) =
  \begin{cases}
    x^2          & \text{se } x \ge 0 \\
    -x           & \text{se } x < 0
  \end{cases}
\end{equation}

\section{Séries}

A soma dos $n$ primeiros inteiros é $\sum_{k=1}^{n} k = \frac{n(n+1)}{2}$.
A série geométrica e a de Basileia:
\begin{align*}
  \sum_{k=0}^{\infty} r^k &= \frac{1}{1 - r}, \qquad |r| < 1 \\
  \sum_{k=1}^{\infty} \frac{1}{k^2} &= \frac{\pi^2}{6}
\end{align*}

\section{Integrais}

$$ \int_{-\infty}^{\infty} e^{-x^2}\,dx = \sqrt{\pi} $$

\begin{gather}
  \int_0^1 x^n\,dx = \frac{1}{n + 1} \\
  \oint_C \vec{F} \cdot d\vec{r} = \iint_S (\nabla \times \vec{F}) \cdot d\vec{S}
\end{gather}

\section{Mecânica e álgebra linear}

A energia de repouso é \(E = mc^2\) e a equação de Schrödinger,
\begin{equation}
  i\hbar \frac{\partial}{\partial t} \Psi(\mathbf{r}, t)
  = \left[ -\frac{\hbar^2}{2m} \nabla^2 + V(\mathbf{r}, t) \right] \Psi(\mathbf{r}, t)
\end{equation}

Uma matriz e o seu determinante:
\[
  A = \begin{pmatrix} a & b \\ c & d \end{pmatrix}, \qquad
  \det A = ad - bc
\]

\begin{multline}
  (x_1 + x_2 + x_3 + x_4)^2 = x_1^2 + x_2^2 + x_3^2 + x_4^2 \\
  + 2(x_1 x_2 + x_1 x_3 + x_1 x_4 + x_2 x_3 + x_2 x_4 + x_3 x_4)
\end{multline}

\section{O que não vira equação}

Custa \$5 e o desconto é de \$2. Um comentário com fórmula fica de fora:
% $a^2 + b^2 = c^2$ (comentário)

Uma equação com erro de sintaxe aparece marcada, sem derrubar as outras: $\frac{1}{$ e, depois dela, $x + y$ normal.

\end{document}
